% !TeX root = ../navier-stokes-manuscript.tex
% ============================================================
% 2.1 Vortex Stretching
% ============================================================

The momentum equation has four pieces:
\[
  \underbrace{\partial_t u}_{\text{local acceleration}}
  +
  \underbrace{(u\cdot\nabla)u}_{\text{nonlinear transport}}
  =
  \underbrace{-\nabla p}_{\text{pressure}}
  +
  \underbrace{\nu\Delta u}_{\text{viscous diffusion}}.
\]
The central regularity problem lies in the competition between nonlinear transport and viscosity.
The hope is simple: viscosity smooths the fluid faster than nonlinear transport can sharpen it.
If that remains true at every scale, the velocity stays smooth.

\subsection{Vortex Stretching}\label{sub:ThePerfectlyStretchingVortex}

A material vortex segment moves with the velocity field that deforms it.
The symmetric part of the velocity gradient,
\[
  S:=\frac12\bigl(\nabla u+(\nabla u)^{\mathsf T}\bigr),
\]
acts on the segment's axial direction \(e\).
If \(L(t)\) is its length and the axial strain has positive rate \(\sigma(t)\), then
\[
  Se=\sigma(t)e,
  \qquad
  \sigma(t)>0,
  \qquad
  \frac{\dot L(t)}{L(t)}
  =
  e\cdot Se
  =
  \sigma(t).
\]
The segment lengthens without acquiring new material.
Its material volume \(\mathcal V\) stays fixed, so with \(A(t):=\mathcal V/L(t)\) and \(r_{\mathrm{eff}}(t):=\sqrt{A(t)/\pi}\), incompressibility gives
\[
  \nabla\cdot u=0,
  \qquad
  \frac{d}{dt}(A L)=0,
  \qquad
  \frac{\dot A}{A}=-\sigma(t),
  \qquad
  \frac{\dot r_{\mathrm{eff}}}{r_{\mathrm{eff}}}
  =
  -\frac{\sigma(t)}{2}.
\]
The axial extension closes the cross-section around the same rotating fluid.
Its rotation is \(\omega:=\nabla\times u\), and along that moving segment it obeys
\[
  \frac{D\omega}{Dt}
  =
  S\omega+\nu\Delta\omega,
  \qquad
  \frac{D}{Dt}
  :=
  \partial_t+u\cdot\nabla,
\]
while alignment with the stretching direction, \(\omega=|\omega|e\), brings the axial strain into its magnitude with the same positive rate:
\[
  S\omega
  =
  \sigma(t)\omega,
  \qquad
  \frac12\frac{D|\omega|^2}{Dt}
  =
  \sigma(t)|\omega|^2
  +
  \nu\,\omega\cdot\Delta\omega.
\]
Diffusion shares this balance, so the rotation strengthens only where the complete right-hand side remains positive.
On such an interval, bounded speed on the fixed-volume segment can leave its kinetic energy finite while vorticity enters the velocity gradient:
\[
  E_{\mathrm{seg}}(t)
  :=
  \frac12\int_{\Omega_{\mathrm{seg}}(t)}|u(x,t)|^2\,dx
  \leq
  \frac12U_0^2\mathcal V
  <
  \infty,
  \qquad
  \left|\frac12\bigl(\nabla u-(\nabla u)^{\mathsf T}\bigr)\right|_{\mathrm F}
  =
  \frac{|\omega|}{\sqrt2}
  \leq
  |\nabla u|_{\mathrm F}.
\]
The gradients rise.

\divider{\NavierStokes{} Scaling}

The narrowing has made width an active variable.
At a fixed time \(t_0\), set \(\ell:=r_{\mathrm{eff}}(t_0)\).
For \(\lambda>1\), the equation's exact finer-scale comparison is
\[
  u_\lambda(x,t)
  :=
  \lambda u(\lambda x,\lambda^2t),
  \qquad
  p_\lambda(x,t)
  :=
  \lambda^2p(\lambda x,\lambda^2t).
\]
At time \(\lambda^{-2}t_0\), its width is \(\lambda^{-1}\ell\), while differentiation gives
\[
\begin{aligned}
  \partial_tu_\lambda(x,t)
  &=
  \lambda^3(\partial_tu)(\lambda x,\lambda^2t),
  \\
  (u_\lambda\cdot\nabla)u_\lambda(x,t)
  &=
  \lambda^3\bigl((u\cdot\nabla)u\bigr)(\lambda x,\lambda^2t),
  \\
  \nabla p_\lambda(x,t)
  &=
  \lambda^3(\nabla p)(\lambda x,\lambda^2t),
  \\
  \nu\Delta u_\lambda(x,t)
  &=
  \lambda^3\nu(\Delta u)(\lambda x,\lambda^2t).
\end{aligned}
\]
The common \(\lambda^3\) leaves viscosity no advantage over nonlinear sharpening at the narrower scale.

\divider{Critical Height}

Severity therefore needs a measure that does not grow merely because the width has changed.
With \(\Lambda:=(-\Delta)^{1/2}\), the homogeneous Sobolev height transforms by
\[
  \norm{u_\lambda(t)}_{\dot H^s}^2
  =
  \lambda^{2s-1}
  \norm{u(\lambda^2t)}_{\dot H^s}^2.
\]
Scale leaves the measurement unchanged only at \(s=\tfrac12\).
At that height set
\[
  H(t)
  :=
  \frac12\norm{u(t)}_{\dot H^{1/2}}^2
  =
  \frac12\norm{\Lambda^{1/2}u(t)}_2^2.
\]
The energy below and the first-derivative energy above move in opposite directions under the same refinement:
\[
  \norm{u_\lambda(t)}_2^2
  =
  \lambda^{-1}\norm{u(\lambda^2t)}_2^2,
  \qquad
  H_\lambda(t)=H(\lambda^2t),
  \qquad
  \norm{\nabla u_\lambda(t)}_2^2
  =
  \lambda\norm{\nabla u(\lambda^2t)}_2^2.
\]
Because the critical height is stationary while its neighboring levels separate, its motion along the original solution must come from the equation.

\divider{Nonlinear Production versus Viscous Dissipation}

Differentiate \(H\) along that solution by applying \(\Lambda^{1/2}\) to the momentum equation and pairing with \(\Lambda^{1/2}u\).
The signed nonlinear contribution is
\[
  \mathcal P_{1/2}(t)
  :=
  -\operatorname{Re}
  \pair
    {\Lambda^{1/2}u(t)}
    {\Lambda^{1/2}\bigl((u(t)\cdot\nabla)u(t)\bigr)}_{L^2}.
\]
Incompressibility removes the pressure pairing, and the Laplacian contributes \(-\nu\norm{\Lambda^{3/2}u}_2^2\), leaving
\[
  H'(t)
  +
  \nu\norm{\Lambda^{3/2}u(t)}_2^2
  =
  \mathcal P_{1/2}(t).
\]
The opening competition is now exact: the height rises precisely when nonlinear production exceeds viscous dissipation.
Refinement acts on both sides of that contest by
\[
  \mathcal P_{1/2}[u_\lambda](t)
  =
  \lambda^2\mathcal P_{1/2}[u](\lambda^2t),
  \qquad
  \nu\norm{\Lambda^{3/2}u_\lambda(t)}_2^2
  =
  \lambda^2\nu\norm{\Lambda^{3/2}u(\lambda^2t)}_2^2.
\]
Their equality of scale persists, but the argument \(\lambda^2t\) places the finer comparison on a time interval shorter by \(\lambda^{-2}\).

\divider{The Heat Clock}

Viscosity already carries this shorter clock.
For the heat equation \(v_t=\nu\Delta v\),
\[
  \widehat v(\xi,t+\delta t)
  =
  e^{-\nu|\xi|^2\delta t}\widehat v(\xi,t).
\]
At width \(r\), the active frequencies have \(|\xi|\simeq r^{-1}\), so order-one diffusion occurs on the time
\[
  \tau_\nu(r)
  :=
  \frac{r^2}{\nu}.
\]
Narrowing the width by \(\lambda^{-1}\) gives
\[
  \tau_\nu(\lambda^{-1}r)
  =
  \lambda^{-2}\tau_\nu(r),
\]
which is the same contraction carried by the time coordinate of the full equation.
The clock belongs to viscosity's response at a given width; the material history must still reach that width.
If its narrowing continues, each new spatial scale arrives with a quadratically shorter duration.

\divider{The Zeno Cascade}

Let the same material segment continue through the geometric widths and their viscous clocks,
\[
  r_n:=2^{-n}r_0,
  \qquad
  \tau_n:=\frac{r_n^2}{\nu},
\]
so these clock lengths contract geometrically and have finite total length:
\[
  \tau_n
  =
  \frac{r_0^2}{\nu}4^{-n},
  \qquad
  \sum_{n=0}^{\infty}\tau_n
  =
  \frac{4r_0^2}{3\nu}
  <
  \infty.
\]
If the segment realizes those widths along its one history, they satisfy
\[
  r_0>r_1>r_2>\cdots\longrightarrow0
\]
at successive passage times
\[
  t_0<t_1<t_2<\cdots\uparrow T_*<\infty,
  \qquad
  t_{n+1}-t_n\asymp\tau_n,
\]
with constants that do not change with \(n\).
The time left after the \(n\)-th passage is then
\[
  T_*-t_n
  \asymp
  \sum_{k=n}^{\infty}\tau_k
  =
  \frac{4r_n^2}{3\nu}.
\]
One material history has crossed infinitely many scales before \(T_*\).
The finite clock does not itself force a loss of regularity; it removes infinite time as the obstruction.

If that same solution loses continuation at \(T_*\) through blow-up in one fixed continuation-controlling Banach space \(B\), none of its finite temporal responses can remain uniformly bounded in \(B\).
Indeed, a bound for any \(\partial_t^N u\) on \([t_0,T_*)\), with \(N\geq1\), integrates over the finite interval to a bound for \(\partial_t^{N-1}u\).
Repeating the integration reaches \(u\) itself and contradicts the assumed blow-up.
Thus blow-up of \(u\) forces loss of uniform \(B\)-control at every finite temporal order.

The \NavierStokes{} equation generates those rows from that same velocity.
Its momentum balance determines \(\partial_tu\); each further time derivative distributes across the nonlinear factors by Leibniz's rule and produces the higher response.
The family \(u,\partial_tu,\partial_t^2u,\ldots\) is the temporal Tower.
